\documentclass[12pt]{article}
\usepackage[utf8]{inputenc}
\usepackage[T1]{fontenc}
\usepackage[margin=1in]{geometry}
\usepackage{setspace}
\usepackage{xcolor}
\usepackage[colorlinks=true, urlcolor=blue!60!black, linkcolor=blue!60!black]{hyperref}
\setlength{\parindent}{2em}
\setlength{\parskip}{1.25em}
\renewcommand{\baselinestretch}{1.0}
\begin{document}
\begin{center}
{\LARGE \textit{Heart of Atlanta Motel, Inc.\ v.\ United States}, 379 U.S.\ 241 (1964)\par}
\vspace{0.8em}
{\large Case Brief --- Four-Page Edition\par}
\end{center}
\vspace{1em}

\section{Memory Jogger}

Congress may use the Commerce Clause to ban racial discrimination by private motels serving interstate travelers, even if its primary motivation is to remedy a moral wrong rather than purely economic harm (\href{https://supreme.justia.com/cases/federal/us/379/241/}{Justia}).

\section{Detailed Case Facts}

Title II of the Civil Rights Act of 1964 prohibited discrimination on the ground of race, color, religion, or national origin in places of public accommodation whose operations ``affect commerce.'' Any inn, hotel, or motel serving transient guests, except an owner-occupied building with no more than five rooms for rent, was declared to affect commerce \textit{per se}. The Heart of Atlanta Motel, a 216-room motel owned by Moreton Rolleston, Jr., sat near Interstates 75 and 85, advertised in national magazines and on more than 50 Georgia billboards, and drew about 75\% of its guests from out of state. It refused to rent rooms to Black guests and intended to continue that policy. On stipulated facts, it challenged Title II under the Commerce Clause, the Fifth Amendment's Due Process and Takings Clauses, and the Thirteenth Amendment.

\section{Procedural History}

The motel sought a declaratory judgment and injunction in the Northern District of Georgia; the United States counterclaimed for enforcement. On July 22, 1964, a three-judge court upheld the relevant provisions and permanently enjoined the motel from refusing Black guests. Justice Black, as Circuit Justice, denied a stay. On direct appeal, argued alongside \textit{Katzenbach v.\ McClung}, the Supreme Court affirmed on December 14, 1964.

\section{Judicial Votes}

Affirmed, \textbf{9--0}.
\begin{itemize}
\item \textbf{Opinion of the Court:} Justice Clark, joined by Chief Justice Warren and Justices Douglas, Harlan, Brennan, Stewart, White, and Goldberg.
\item \textbf{Concurrences:} Justice Black (in the judgment and holding); Justice Douglas; Justice Goldberg.
\item \textbf{Dissents:} None.
\end{itemize}

\section{Holding}

Congress may prohibit racial discrimination by public accommodations serving interstate travelers when it has a rational basis for finding that the discrimination affects interstate commerce and the means chosen are reasonable and appropriate. The interstate movement of persons is commerce, and Congress may regulate its local incidents even while addressing a moral and social wrong. Title II violates neither the Fifth nor the Thirteenth Amendment. The Court declined to decide whether \S~5 of the Fourteenth Amendment also supported Title II, and its remark that the motel would likely suffer no long-run loss is dicta.

\section{Analysis of Opinions}

\textbf{Majority (Clark).} The \textit{Civil Rights Cases} (1883) did not control: Title II was limited to enterprises substantially related to interstate commerce, the 1883 Court had not considered the Commerce Clause, and interstate travel had expanded enormously. The hearings were ``replete with evidence'' that lodging discrimination burdened Black travelers, forcing them to rely on special guidebooks, drive long distances, or forgo travel. From \textit{Gibbons v.\ Ogden} through \textit{Hoke} and \textit{Caminetti}, the movement of persons is commerce, and Congress may remove local obstructions to it. Thirty-two States already barred such discrimination, so there was no due process problem; economic loss from valid regulation is no taking; and Title II codified the old innkeeper rule, bearing no resemblance to slavery.

\textbf{Black (concurring).} Title II rests on the Commerce and Necessary and Proper Clauses. Not ``every remote, possible, speculative effect on commerce'' suffices, and the question is ultimately judicial; but under \textit{Wickard}'s aggregation principle, widespread exclusion severely burdened interstate travel.

\textbf{Douglas (concurring).} The right to be free of discrimination ``occupies a more protected position in our constitutional system than does the movement of cattle, fruit, steel and coal across state lines.'' He would rest Title II on \S~5 and \textit{Shelley v.\ Kraemer}, for a ``more settling effect.''

\textbf{Goldberg (concurring).} The Act's purpose is human dignity; Congress had authority under both the Commerce Clause and \S~5.

\section{Examples: Future Applications}

\begin{enumerate}
\item \textbf{Same-Side: Vacation Rentals.} A ban on disability or familial-status discrimination by commercial vacation-rental hosts with more than five units, backed by hearings, regulates lodging tied to interstate travel.
\item \textbf{Same-Side: Highway EV Chargers.} Barring charging stations along highways from refusing drivers by registration or protected traits regulates local incidents of interstate travel.
\item \textbf{Opposite-Side: Gender-Motivated Violence.} A civil remedy for local, non-commercial violence fails under \textit{Morrison} (2000); unlike motel operations, the activity is non-economic.
\item \textbf{Opposite-Side: Guns Near Schools.} Banning gun possession near a school without a jurisdictional nexus fails under \textit{Lopez} (1995).
\item \textbf{Fence-Sitter: Repealing ``Mrs.\ Murphy.''} Covering a remote two-bedroom farmhouse with only local advertising and in-state guests tests aggregation against Black's warning about ``remote, possible, speculative'' effects.
\end{enumerate}

\section{Critique}

\textbf{Originalist and libertarian.} Randy Barnett, Richard Epstein, and Justice Thomas question the ``substantial effects'' doctrine, arguing that commerce originally meant trade in goods and that deferential review turns the clause into a police power; Epstein adds that Title II overrides association and property rights beyond the common-law innkeeper rationale.

\textbf{Progressive.} Many scholars applaud the result but echo Douglas and Goldberg: resting Title II on commerce and leaving the \textit{Civil Rights Cases} intact framed equality in commercial terms, leaving Congress without a Fourteenth Amendment foothold once \textit{Morrison} and \textit{City of Boerne} narrowed federal power.

\section{Key Quotations}

\begin{quote}
``Thus, the power of Congress to promote interstate commerce also includes the power to regulate the local incidents thereof, including local activities in both the States of origin and destination, which might have a substantial and harmful effect upon that commerce.''
\end{quote}
\noindent --- Justice Clark, for the Court, 379 U.S.\ at 258.

\begin{quote}
``In short, the determinative test of the exercise of power by the Congress under the Commerce Clause is simply whether the activity sought to be regulated is `commerce which concerns more States than one' and has a real and substantial relation to the national interest.''
\end{quote}
\noindent --- Justice Clark, for the Court, 379 U.S.\ at 255.

\begin{quote}
``It was this burden which empowered Congress to enact appropriate legislation, and, given this basis for the exercise of its power, Congress was not restricted by the fact that the particular obstruction to interstate commerce with which it was dealing was also deemed a moral and social wrong.''
\end{quote}
\noindent --- Justice Clark, for the Court, 379 U.S.\ at 257.

\end{document}
